\documentclass[UTF8,fontset=none,11pt,a4paper]{ctexart}
\usepackage{ipara-handbook}
\handbooksetup
  {DS-03 串与模式匹配}
  {408 · 数据结构 DS-03}
  {串 · 前缀 · 失配 · 复用}
\begin{document}
\handbooktitle
  {串与模式匹配}
  {失配之后，怎样复用已经确认的前缀信息}
  {408 · 数据结构专题 DS-03}

\begin{corebox}{母问题：失配之后，哪些信息真的已经知道？}
在一条较长主串中查找较短模式串时，每次失配之前其实已经证明了一件事：主串刚刚匹配成功的那一段后缀，与模式串开头的一段前缀完全相同。KMP 不把这条事实丢掉，而是继续问：\kw{这段已经确认相等的后缀中，最长有多少字符还能直接作为模式串前缀继续使用？}

后文先用一个具体失配例子把这个问题跑通，再引入主串 $T$、模式串 $P$、下标和前缀函数。这样所有符号都有具体对象可指，不先靠抽象公式定义算法。
\end{corebox}

\begin{methodbox}{本册定位与归属边界}
\begin{itemize}
\item \kw{本册负责：}串的基本对象、朴素匹配、前缀函数、失配回退信息与 KMP 主循环。
\item \kw{调用：}数据结构总图的表示—操作—不变量—成本语言；底层字符序列只借用一般线性存储，不重新定义 DS01。
\item \kw{输出：}向算法专题提供线性时间模式匹配的可调用接口。
\item \kw{边界：}完整确定有限自动机、多模式匹配与文本索引属于后续扩展，不在本册展开。
\end{itemize}
\end{methodbox}

\tableofcontents\newpage

\section{先从一次真实失配看懂 KMP 在做什么}
先固定两个对象：$T$ 表示\kw{主串}，即被搜索的较长字符串；$P$ 表示\kw{模式串}，即希望在主串中找到的较短字符串。模式匹配的输出契约是：返回 $P$ 在 $T$ 中第一次出现的起始位置，若不存在则返回 `npos`。本册代码另约定空模式返回 0。

\begin{examplebox}{同一个主串字符为什么不用回头重比}
取
\[
T=\texttt{abababaca},\qquad P=\texttt{ababaca}.
\]
采用零起始下标。前 5 个字符已经匹配成功，此时准备比较的是主串字符 $T[5]=\texttt{b}$ 与模式字符 $P[5]=\texttt{c}$，因此在 $i=5,j=5$ 处失配。

失配之前已经知道
\[
T[0\ldots4]=P[0\ldots4]=\texttt{ababa}.
\]
而 `ababa' 的末尾 `aba' 与开头 `aba' 相同，所以最后 3 个已经匹配的字符无需重新比较。模式状态可以直接从 $j=5$ 缩短到 $j=3$，主串位置仍保持 $i=5$。下一次真正比较的是
\[
T[5]=\texttt{b}\qquad\text{与}\qquad P[3]=\texttt{b},
\]
这次比较成功。

因此 KMP 的关键不是先问“模式串右移几格”，而是先问：\kw{已经匹配成功的文本后缀中，最长有多少字符也等于模式前缀？}
\end{examplebox}

现在再抽象。若当前准备比较 $T[i]$ 与 $P[j]$，并且前面已有 $j$ 个字符匹配成功，则保持
\[
T[i-j\ldots i-1]=P[0\ldots j-1].
\]
所以 $i$ 表示当前主串位置，$j$ 表示当前模式位置；在这个零起始状态模型里，$j$ 同时也是已经匹配成功的字符数。失配后，KMP 只缩短 $j$，尽量保留上式已经证明过的相等关系。

\subsection{朴素匹配为什么会重复工作}
朴素匹配枚举起点 $s=0,1,\ldots,\lvert T\rvert-\lvert P\rvert$，在每个起点从左到右比较。若在第 $j$ 个字符失配，已经比较过的前缀信息只用于判定当前起点失败，算法不会把它带到下一起点。最坏比较次数可以达到 $\Theta((n-m+1)m)$，其中 $n=\lvert T\rvert,m=\lvert P\rvert$。

例如文本 $T=\texttt{aaaaab}$、模式 $P=\texttt{aaab}$ 时，多个起点都会先匹配若干个 `a' 再失配。朴素算法反复重新确认这些 `a'；KMP 则把模式自身的重复结构预先记录下来，让失配后的新状态直接继承已经确认的信息。

\subsection{朴素匹配的两个指针与回退公式}
朴素匹配常用文本指针 $i$ 和模式指针 $j$。当 $T[i]=P[j]$ 时同时前进；失配且 $j>0$ 时，当前起点已经失败，模式回到 $0$，文本应回到“该起点后的下一个字符”：
\[
i\leftarrow i-j+1,\qquad j\leftarrow0.
\]
如果代码采用“比较失败后循环末尾再统一执行 $i++$”，就会把同一语义写成 $i\leftarrow i-j+2$。两个公式不是不同算法，而是指针更新时机不同；判断是否正确要看下一轮比较的文本位置，而不是死记加一的数值。

朴素匹配的不变量是：进入每个起点时，$i$ 指向该起点，$j=0$；循环结束时，已经比较的 $T[s..i-1]$ 与 $P[0..j-1]$ 相等，或已找到第一个不等字符。它的最坏界来自很多起点重复确认同一主串片段，而不是来自“回退”这个词本身。用朴素匹配作为基线，可以明确 KMP 真正删除的是“重新建立匹配前缀”的重复工作。

\section{前缀函数：模式自身到底有哪些信息可以复用}
先固定术语。前缀必须从字符串第一个字符开始，后缀必须在最后一个字符结束；真前缀和真后缀不能等于整个字符串。若一个字符串同时是真前缀和真后缀，本册称它为\kw{边界}。例如 `ababa' 的真前缀有 `a, ab, aba, abab'，真后缀有 `a, ba, aba, baba'，共同部分为 `a, aba'，所以最长边界是 `aba'，长度为 3。

为了不和匹配阶段的主串下标 $i$ 混淆，这一节统一用 $r$ 表示“当前正在构造前缀函数的模式串位置”。对模式前缀 $P[0..r]$，定义 $\pi[r]$ 为其最长边界的\kw{长度}：
\[
\pi[r]=\max\{k:\;0\le k<r+1,\;P[0..k-1]=P[r-k+1..r]\}.
\]
计算 $\pi[r]$ 时，用 $k$ 表示“当前候选边界的长度”。先令 $k=\pi[r-1]$，意思是“先尝试把上一位置的最长边界再延长一个字符”。若 $P[r]\ne P[k]$，长度 $k+1$ 的边界不可能成立。此时更短但仍可能成立的候选，必须是长度 $k$ 这段前缀自己的边界，所以直接令
\[
k\leftarrow\pi[k-1]
\]
继续尝试。若最终有 $P[r]=P[k]$，则令 $k\leftarrow k+1$；若所有非空候选都失败，则保留 $k=0$。这里特意不用字母 $b$ 表示候选长度，避免与模式串中实际字符 `b' 混淆。

因此构造不变量非常具体：处理完 $P[0..r]$ 后，数组第 $r$ 项保存的就是该前缀最长边界的长度。每次回退都严格缩短候选长度，不会在已经排除的候选上反复比较。

\subsection{把 \texorpdfstring{$\pi$}{pi} 真正算一遍：\code{ababaca}}
继续沿用上面的模式串 $P=\texttt{ababaca}$。逐个加入新字符，并始终问“当前最长边界能否延长；不能时，这个边界自己的最长边界是谁”：
\begin{center}\small
\begin{tabularx}{\linewidth}{C{0.8cm}C{1.1cm}L{2.7cm}Y C{1.1cm}}
\toprule
$r$ & $P[r]$ & 旧候选长度 $k$ & 比较、回退与可复用事实 & $\pi[r]$\\
\midrule
0 & a & 无 & 单字符没有非空真边界 & 0\\
1 & b & 0 & $b\ne a$，无法延长 & 0\\
2 & a & 0 & $a=a$，空边界延长为 `a' & 1\\
3 & b & 1 & $b=P[1]$，`a' 延长为 `ab' & 2\\
4 & a & 2 & $a=P[2]$，`ab' 延长为 `aba' & 3\\
5 & c & 3 & $c\ne P[3]$：$3\to1\to0$，仍不等 & 0\\
6 & a & 0 & $a=P[0]$，重新得到 `a' & 1\\
\bottomrule
\end{tabularx}
\end{center}
因此 $\pi=[0,0,1,2,3,0,1]$。第 5 项的候选长度依次从 3 缩到 1、再缩到 0，是因为长度 3 的候选已经失败；后续仍可能成立的更短候选，只可能来自这段长度 3 前缀自己的边界链。

\section{KMP 主循环：同一个文本字符为什么可以继续比较}
为了让数学符号与代码变量分开，下面用 $q$ 表示“已经匹配成功的字符数”；C++ 实现中的变量 `matched' 保存的就是这个 $q$。在准备处理当前主串字符 $T[i]$ 之前，保持不变量
\[
T[i-q\ldots i-1]=P[0\ldots q-1].
\]
若 $T[i]\ne P[q]$ 且 $q>0$，当前最长状态无法继续，但上式已经证明的文本后缀并没有消失，因此令
\[
q\leftarrow\pi[q-1]
\]
改试更短的边界。注意此时\kw{同一个 $T[i]$ 仍未被消费}，所以代码中的 `while' 会继续比较，而不是移动主串下标。只有找到可以接上的模式位置，或所有非空候选都失败以后，外层扫描才进入下一个主串字符。

若当前字符能够接上，则 $q$ 增加 1；当 $q=m$ 时，模式全部匹配完成，匹配起点为 $i+1-m$。对应到代码，就是 `matched++' 与 `matched == pattern.size()'。

\begin{examplebox}{一次完整失配：文本不回头，模式状态缩短}
继续使用 $T=\texttt{abababaca}$、$P=\texttt{ababaca}$。当 $i=5$ 时，前面已经匹配成功 5 个字符，因此 $q=5$。当前真正比较的是主串字符 $T[5]=\texttt{b}$ 与模式字符 $P[5]=\texttt{c}$，二者失配。

由前缀函数可知 $\pi[4]=3$，所以把已匹配长度从 5 缩短到 3。主串位置 $i$ 不变，下一次真正比较的是同一个 $T[5]=\texttt{b}$ 与 $P[3]=\texttt{b}$，这次匹配成功，于是 $q$ 增加到 4。整个过程中没有重新比较已经证实相等的 `aba'。
\end{examplebox}

\begin{center}\small
\begin{tabularx}{\linewidth}{L{2.3cm}YY}
\toprule
状态 & 迁移 & 必须保持的事实\\
\midrule
$q=0$ & 比较当前主串字符与 $P[0]$ & 没有可复用的非空模式前缀\\
$q>0$ 且失配 & 沿前缀函数链缩短 $q$ & 不移动主串下标，保留已证实的后缀\\
匹配后 & $q\leftarrow q+1$ & 当前状态代表长度为 $q$ 的相等前缀\\
$q=m$ & 返回 $i+1-m$ & 所有模式字符已连续匹配\\
\bottomrule
\end{tabularx}
\end{center}

\subsection{KMP 的线性时间证明：把回退次数当作势能}
主串下标 $i$ 从不减小；已匹配长度 $q$ 每次成功比较最多增加 1，每次失配回退都严格跳到更短的边界。一次失配可能沿前缀函数链回退多次，但每次回退都消耗此前积累的 $q$。令势能 $\Phi=q$：成功比较最多让势能增加 1，回退则严格降低势能，所以整个主串扫描中的回退总量被此前的增长次数所限制。于是扫描时间为 $O(n)$，再加模式前缀函数的 $O(m)$ 构造，总时间为 $O(n+m)$。

\subsection{教材中的 \code{next} 与 \code{nextval}：先固定状态语义}
不同教材会把失配回退信息编码成不同数组；数组首项不是算法本体。为了和 408 常见手算口径建立一个可检验接口，本册额外固定一套\kw{零起始下标 + $-1$ 哨兵}的教材编码。设当前正在比较 $T[i]$ 与 $P[j]$，并且失配前已有
\[
T[i-j\ldots i-1]=P[0\ldots j-1].
\]
于是 $j$ 既是当前待比较的模式下标，也是已经匹配成功的字符数。若 $j>0$，失配后仍可能复用的最长状态只取决于已匹配部分 $P[0\ldots j-1]$，因此
\[
\operatorname{next}[0]=-1,\qquad
\operatorname{next}[j]=\pi[j-1]\quad(j>0).
\]
这就是“前缀函数右移一位、首位补 $-1$”的真正原因：\code{next[j]} 回答的是“第 $j$ 位失配以后怎么办”，所以读取的是失配前那 $j$ 个字符的最长边界，而不是 $P[0\ldots j]$ 的最长边界。

若
\[
P[j]=P[\operatorname{next}[j]],
\]
则已经知道 $T[i]\ne P[j]$ 时，回退后再拿同一个 $T[i]$ 与相同字符比较必然再次失败。于是可继续跳过这个无效状态：
\[
\operatorname{nextval}[j]=
\begin{cases}
-1,&j=0,\\
\operatorname{nextval}[\operatorname{next}[j]],&j>0\text{ 且 }P[j]=P[\operatorname{next}[j]],\\
\operatorname{next}[j],&j>0\text{ 且 }P[j]\ne P[\operatorname{next}[j]].
\end{cases}
\]
这里 $-1$ 只是哨兵状态，\kw{不是长度为 $-1$ 的前缀}。因此把 \code{nextval[j]} 直接解释成“新对齐前缀长度”只在其非负时成立；更稳定的统一语言是“失配后的下一模式状态编码”。

\begin{examplebox}{从同一个 prefix 生成教材数组：\code{aabaab}}
对 $P=\code{aabaab}$，
\[
\pi=[0,1,0,1,2,3],
\]
故
\[
\operatorname{next}=[-1,0,1,0,1,2],
\qquad
\operatorname{nextval}=[-1,-1,1,-1,-1,1].
\]
例如 $j=4$ 失配时，\code{nextval[4]=-1} 表示没有值得用同一文本字符继续尝试的非负模式状态；标准哨兵主循环随后执行 $i\mathrel{+}=1,j\mathrel{+}=1$，下一次真正比较从模式下标 0 开始。
\end{examplebox}

\subsection{模式串为什么右移 \texorpdfstring{$j-j'$}{j-j'}：从起点坐标生成公式}
设失配前模式起点为
\[
s=i-j.
\]
若失配后状态变为非负的 $j'$，主串位置 $i$ 不变，新模式起点为
\[
s'=i-j'.
\]
因此模式右移量不是 \code{next} 值本身，而是
\[
\Delta=s'-s=j-j'.
\]
在本节的修正数组口径下取 $j'=\operatorname{nextval}[j]$，得到
\[
\boxed{\Delta=j-\operatorname{nextval}[j]}.
\]
当 $j'=-1$ 时，下一步哨兵迁移使真正的新起点成为 $i+1$，故
\[
(i+1)-(i-j)=j+1=j-(-1),
\]
同一公式仍成立。这个边界推导很重要：公式成立来自状态机的下一步行为，而不是把 $-1$ 误当成某种前缀长度。

\subsection{输出契约：首个匹配与全部匹配不是同一个接受动作}
匹配接口还要区分“返回首个位置”和“枚举全部出现”。若允许重叠匹配，成功一次后不能把已匹配长度直接清零，而应令
\[
q\leftarrow\pi[q-1],
\]
继续扫描下一个主串字符；例如模式 `aba' 在 `ababa' 中出现于 0 和 2。若只返回首个位置，接受状态即可立即结束。接口契约改变，状态复位策略随之改变。

\section{代码：把两种扫描策略放在同一契约下}
完整 C++17 实现位于 \codepath{code/ds03_string_matching.hpp}，测试位于 \codepath{code/ds03_string_matching_test.cpp}；正文代码块直接从真实源文件抽取。

\subsection{朴素基线：每个起点重新建立 matched}
\lstinputlisting[
  language=C++,firstline=10,lastline=27,firstnumber=1,numbers=left,
  caption={朴素匹配的起点枚举与局部比较}
]{30_408/10_数据结构/03_串与模式匹配/code/ds03_string_matching.hpp}

朴素版本的重要性在于提供复杂度基线：它没有错误，只是没有利用跨起点的结构信息。任何声称 KMP 更快的解释，都应能指出它省掉了哪类重复比较。

\subsection{前缀函数：失配时只沿边界链回退}
\lstinputlisting[
  language=C++,firstline=31,lastline=46,firstnumber=1,numbers=left,
  caption={前缀函数的边界回退构造}
]{30_408/10_数据结构/03_串与模式匹配/code/ds03_string_matching.hpp}

代码变量 `border' 只表示当前候选边界的长度。每次回退都跳到更短边界，因此不会形成无限循环；若新字符与候选前缀的下一字符相等，再增长一次。`prefix' 数组保存的就是每个模式前缀的最长边界长度。

\subsection{KMP：失配不倒退主串}
\lstinputlisting[
  language=C++,firstline=50,lastline=68,firstnumber=1,numbers=left,
  caption={KMP 主循环与完整模式接受}
]{30_408/10_数据结构/03_串与模式匹配/code/ds03_string_matching.hpp}

这段代码最容易误读的是 `while'。它不是“重新开始匹配”，而是\kw{在不移动当前主串字符的前提下，连续缩短可复用模式状态}。代码变量 `matched' 保存前文数学状态 $q$。失配时，代码读取 `prefix[matched - 1]'；由于该数组项保存的正是 $\pi[q-1]$，所以它和前文的数学回退完全一致。

主串循环变量 `i` 单调递增；模式状态 `matched` 才可能沿前缀函数链下降。于是预处理 $O(m)$ 加扫描 $O(n)$，总时间 $O(n+m)$，额外空间 $O(m)$。这不是“next 公式自动带来线性”，而是由文本扫描不倒退、模式状态回退严格变短共同保证。

\subsection{教材编码：从前缀函数确定性生成 \code{next/nextval}}
\lstinputlisting[
  language=C++,firstline=70,lastline=106,firstnumber=1,numbers=left,
  caption={零起始、$-1$ 哨兵的 \code{next} 与 \code{nextval} 生成}
]{30_408/10_数据结构/03_串与模式匹配/code/ds03_string_matching.hpp}

这两个函数故意建立在 \code{prefix\_function} 之上：先固定唯一数学对象，再生成考试编码，避免让 \code{next} 与 \code{nextval} 变成两套平行理论。以 \code{aabaab} 的 $j=4$ 为例，\code{next[4]=1}，但 $P[4]=P[1]=\code{a}$；如果当前文本字符已经与 $P[4]$ 的 `a' 失配，再退到同样为 `a' 的 $P[1]$ 比较必然仍失败，所以 \code{nextval[4]} 继续取 \code{nextval[1]=-1}。

\subsection{教材版主循环：让代码和手算使用同一组 \texorpdfstring{$(i,j)$}{(i,j)} 状态}
\lstinputlisting[
  language=C++,firstline=108,lastline=134,firstnumber=1,numbers=left,
  caption={零起始、$-1$ 哨兵的教材版 KMP 主循环}
]{30_408/10_数据结构/03_串与模式匹配/code/ds03_string_matching.hpp}

这里 $i$ 指向当前文本字符，$j$ 指向当前模式字符；在普通失配分支中只有
\[
j\leftarrow\operatorname{next}[j],
\]
没有 $i\mathrel{+}=1$，所以同一个 $T[i]$ 会继续与回退后的 $P[j]$ 比较。若最终回退到 $j=-1$，下一轮进入哨兵分支，同时令 $i\mathrel{+}=1,j\mathrel{+}=1$，把模式首字符移到下一个文本起点。于是手算中的“$j$ 回退、$i$ 不动”与代码是一一对应的状态转移，而不是口号。

前缀函数版和教材版的对应关系可以压成四项：代码变量 `matched' 与教材状态 $j$ 都表示已经匹配成功的字符数；代码查 `prefix[matched - 1]'，教材查 `next[j]'，两者都在寻找失配后的新模式状态；普通失配时当前主串字符都不被消费；代码中 `matched == pattern.size()' 与教材中的 $j=m$ 都表示模式全部匹配成功。

\subsection{修正数组版：主循环不变，只替换失配表}
\lstinputlisting[
  language=C++,firstline=136,lastline=158,firstnumber=1,numbers=left,
  caption={使用 \code{nextval} 的教材版 KMP 主循环}
]{30_408/10_数据结构/03_串与模式匹配/code/ds03_string_matching.hpp}

`nextval' 版本没有改写 KMP 主循环的基本语义。匹配成功仍然同时推进主串与模式位置；普通失配仍然保持当前主串字符，只修改模式状态。区别只有失配表由 `next' 换成 `nextval'：后者提前越过模式字符相同、已经能够判定必败的中间状态。因此它改变的是某些输入上的字符比较次数，不改变匹配结果，也不改变 $O(n+m)$ 的渐近上界。

\subsection{测试：重复前后缀、教材编码与三种 KMP 等价性}
\lstinputlisting[
  language=C++,firstline=16,lastline=65,firstnumber=1,numbers=left,
  caption={边界输入、三种 KMP 对拍、前缀函数与教材数组}
]{30_408/10_数据结构/03_串与模式匹配/code/ds03_string_matching_test.cpp}

测试要求朴素匹配、前缀函数版 KMP、普通 `next' 版 KMP 和 `nextval' 版 KMP 在同一批主串/模式串上返回相同的首个位置，同时直接检查前缀函数与教材数组的数值，并覆盖空模式、模式过长、完全无匹配和失配后仍有答案的情形。

\lstinputlisting[
  language=C++,firstline=67,lastline=107,firstnumber=1,numbers=left,
  caption={字符比较次数与 \code{nextval} 优化的可执行验证}
]{30_408/10_数据结构/03_串与模式匹配/code/ds03_string_matching_test.cpp}

这里特意把“字符比较”定义为真正执行一次主串字符与模式字符的比较；进入 $j=-1$ 的哨兵分支只更新状态，不计字符比较。测试锁定 2019 年真题主串 `abaabaabcabaabc'、模式串 `abaabc' 的比较次数为 10；同时用主串 `aaaacaaaaab'、模式串 `aaaaab' 验证普通 `next' 需要 15 次字符比较，而 `nextval' 只需要 11 次，且二者匹配结果相同。

测试还枚举二字母表上长度不超过 6 的全部主串、长度不超过 4 的全部模式串，逐对验证三种 KMP 与朴素匹配的结果一致。穷举不是正确性证明，但会系统攻击空模式、重叠匹配、周期串、长边界链与无匹配状态。

\begin{lstlisting}[language=bash]
clang++ -std=c++17 -Wall -Wextra -Werror -pedantic \\
  code/ds03_string_matching_test.cpp -o /tmp/ds03_test
/tmp/ds03_test
\end{lstlisting}

\section{复杂度与边界}
\begin{center}\small
\begin{tabularx}{\linewidth}{L{2.5cm}C{1.8cm}C{1.8cm}YY}
\toprule
任务 & 朴素匹配 & KMP & 关键前提 & 失败边界\\
\midrule
预处理模式 & 无 & $O(m)$ & 前缀函数语义固定 & 空模式\\
扫描文本 & $O((n-m+1)m)$ 最坏 & $O(n)$ & 主串下标单调前进 & $m>n$\\
总额外空间 & $O(1)$ & $O(m)$ & 只保存模式信息 & 大模式内存预算\\
返回位置 & 首个合法起点 & 首个合法起点 & 接受后立即返回 & 无匹配返回 npos\\
\bottomrule
\end{tabularx}
\end{center}

不同教材的 `next` 可能使用“失配时跳到的下标”或“最长边界长度”两种编码。做题时先写语义，再把题目下标转换到该语义；不能拿另一版本的初值和边界直接替换。

\subsection{串对象与存储边界}
串是字符受限的线性表，但题面中的“字符个数”“存储长度”“子串长度”仍需分别解释。长度为 $n$ 的串有 $n(n+1)/2$ 个非空连续子串；若把空串也算入则再加 1。相同内容出现在不同位置时，在“子串个数”问题中是否去重必须由题设声明。

顺序串可以用固定数组、动态数组或长度字段加字符区表示；链式串可以把字符分块存储。表示改变的是随机访问、连接、扩容和额外空间成本，不改变子串与模式匹配的逻辑定义。若题目给出“长度不计终止符”，C 风格终止字符只是实现边界；若问存储容量，则必须把终止符、块内空闲字符或长度字段的约定单独计入。

\begin{methodbox}{从长度语义翻译到教材 next}
本册的前缀函数只维护 $\pi[r]=$“$P[0..r]$ 的最长边界长度”。若教材把 `next[j]' 定义为“模式状态 $j$ 失配后跳到的状态”，零基长度状态下对应 $\pi[j-1]$（$j>0$）。若教材使用 $-1$ 哨兵或一基下标，先画出“已经匹配多少个字符”，再平移下标；数组首项长什么样不是算法本体。
\end{methodbox}

\subsection{从单模式扫描到全文搜索：索引改变了成本归属}
对多篇文档做查询时，逐篇运行朴素匹配或 KMP 的成本为“文档总字符数乘以查询次数”。KMP 只减少单篇扫描的失配浪费，并没有把文档集合变成可按词定位的结构。若查询重复出现，可以离线建立倒排索引：
\[
\text{document}\to\text{token/term}\to\text{posting list of document IDs}.
\]
此时查询阶段从扫描正文转为查索引，更新文档则要承担分词、词项维护和倒排表更新成本。倒排索引与 KMP 都是“预处理换查询”，但保存的信息不同：KMP 保存模式自身的边界结构，倒排索引保存语料中的出现位置。不能因为两者都使用预处理，就把一个当成另一个的实现。

全文搜索还会引入大小写、分词、规范化、通配符和多模式查询等语义问题；这些属于扩展内容。408 题目若只给一个主串和一个模式串，优先调用本册的单模式契约，不凭“工业搜索”背景增加不存在的预处理。

\section{概念边界与做题控制协议}
\begin{center}\small
\begin{tabularx}{\linewidth}{L{1.9cm}C{0.35cm}L{1.9cm}YY}
\toprule
概念 A & $\neq$ & 概念 B & 真正区别 & 混淆后果\\
\midrule
串 & $\neq$ & 模式 & 被搜索的文本对象 vs 约束模板 & 方向写反\\
边界 & $\neq$ & 任意重复片段 & 必须同时是前缀和后缀 & 错误回退位置\\
前缀函数 & $\neq$ & 匹配答案 & 模式内部信息 vs 文本中的出现位置 & 把预处理当搜索\\
失配回退 & $\neq$ & 主串回退 & 只缩短模式状态 & 复杂度退化\\
线性时间 & $\neq$ & 每次比较 $O(1)$ & 总比较次数还需证明单调性 & 复杂度口号化\\
\bottomrule
\end{tabularx}
\end{center}

遇到匹配题，按“固定返回契约 → 写朴素基线 → 找重复前后缀 → 固定 prefix/next 语义 → 画一次失配回退 → 再报告复杂度”的顺序处理。忘掉公式后，只需复原：\kw{已经匹配了什么？失配时哪些后缀仍可能成为模式前缀？模式状态能否缩短而不重扫文本？}

\section{全科坐标与来源处理}
\begin{corebox}{全科坐标：前缀 / 失配 / 复用}
DS03 把线性序列上的“匹配失败”转化为模式内部的边界结构；它接收 DS01 的顺序访问成本语言，并向算法专题输出 KMP 的线性扫描接口。自动机、多模式和文本索引属于后续扩展，不改变本册的唯一知识归属。
\end{corebox}

本册吸收归档总册对 KMP 的核心表述“已匹配前缀压缩为失败状态”，并将 next 下标差异明确为表示约定问题。代码和测试只验证本册固定的 prefix-length 语义，避免把不同教材的数组编码混作第二套稳定真相。
\end{document}
